\subsection{معرفی و انگیزه}
این ماژول در چارچوب پژوهش حاضر پیاده‌سازی و ارزیابی شده و فاقد مقاله مستقل جداگانه در ادبیات تحقیق است.

مکانیزم توجه خطی گیت‌دار \lr{(Gated Linear Attention - GLA)} به عنوان پلی نوآورانه میان مدل‌های توجه استاندارد ترنسفورمر با پیچیدگی محاسباتی و حافظه‌ای مرتبه دوم $\mathcal{O}(N^2)$ و مدل‌های حالت فضای خطی \lr{(SSM / RNN)} با پیچیدگی $\mathcal{O}(N)$ توسعه یافته است. در الگوهای سنتی خودتوجهی، هزینه پردازش و ردپای حافظه با افزایش طول دنباله ورودی به صورت مربعی منفجر شده و مانع از مقیاس‌پذیری کارآمد در پردازش متون و اسناد بسیار طولانی می‌گردد. اگرچه تقریب‌های پیشین توجه خطی با استفاده از ویژگی‌های کرنلی ضرب‌پذیر پیچیدگی را به $\mathcal{O}(N)$ رساندند، اما فقدان مکانیزم تضعیف وابسته به داده منجر به فروپاشی ظرفیت حافظه تداعی‌گر و فراموشی شدید اطلاعات متنی دوربرد در مقایسه با تبدیل غیرخطی بیشینه هموار \lr{(Softmax)} می‌شد.

ایده محوری در ماژول \lr{GatedLinearAttention}، تلفیق بازگشت‌پذیری خطی چندسره با گیت‌های زوال وابسته به ورودی \lr{(Data-Dependent Decay Gating)} و گیتینگ ضربی خروجی بر پایه تابع فعال‌ساز \lr{SiLU} است. این معماری به هر سر توجه اجازه می‌دهد به صورت کامپکت و پویا، نرخ نگهداری و میرایی اطلاعات در طول زمان را تنظیم نموده و پایداری ماتریس حافظه را تضمین نماید. همچنین به کارگیری نرمال‌سازی مخرج کسری مانع از سرریز یا ناپایداری‌های مقیاسی در دنباله‌های طولانی می‌گردد.
